% =====================================================================
% CAPITOLO 2 — BACKGROUND
% =====================================================================
% Scopo (linee guida del prof): mettere un lettore non specialista
% (informatica magistrale) in condizione di apprezzare i contributi.
% Spiegare SOLO i concetti che riappaiono nei capitoli polpa.
% NON una rassegna esaustiva: una raccolta categorizzata costruita
% mentre si scrive la polpa.
%
% Ordine di scrittura: SECONDO (subito dopo la polpa).
% Lunghezza indicativa: ~25-30 pagine.
% =====================================================================

\chapter{Background}
\label{chap:background}

% TODO — Aprire il capitolo con un breve "manifesto" dello stile:
% spiegazione concisa, ogni sezione ancorata a 1-2 riferimenti
% bibliografici per l'approfondimento, niente di superfluo.

% =====================================================================
\section{Dai modelli di linguaggio ai Transformer}
\label{sec:bg-transformer}
% =====================================================================
% TODO — Richiamo storico SCURNO: n-gram -> RNN/LSTM -> attenzione ->
% Transformer \cite{vaswani2017attention}. Perche' il Transformer
% ha vinto: parallelizzazione, dipendenze a lungo raggio, scalabilita'.
% Distinzione encoder/decoder/encoder-decoder; dominanza decoder-only
% (GPT, LLaMA, Gemini, Claude).

% =====================================================================
\section{Large Language Model: capacita', limiti, allineamento}
\label{sec:bg-llm}
% =====================================================================
% TODO — Pre-training autoregressivo, capacita' emergenti, leggi di scala
% \cite{brown2020gpt3}. Instruction tuning + RLHF \cite{ouyang2022instructgpt}
% come allineamento al comportamento desiderato.
% Allucinazioni come limite strutturale \cite{xu2024hallucination}.
% Knowledge cutoff e conseguenze pratiche in azienda.

% =====================================================================
\section{Prompt engineering, function calling, tool use}
\label{sec:bg-prompting}
% =====================================================================
% TODO — Zero-shot, few-shot, Chain-of-Thought \cite{wei2022cot}.
% Function calling: schema, agentic loop, tool nidificati.
% E' il meccanismo su cui si fondano gli agenti Text-to-SQL e il router
% di Orion (output strutturato).

% =====================================================================
\section{Grounding: ancorare la generazione a fonti esterne}
\label{sec:bg-grounding}
% =====================================================================
% TODO — Il grounding come contromisura alle allucinazioni
% \cite{guu2020realm}. Differenza concettuale tra grounding (intervento
% a inferenza) e fine-tuning (modifica dei pesi).

% [FIGURA 2.1] Tre meccanismi di grounding affiancati: Web grounding,
% RAG, Text-to-SQL — uno schema che mostri input, fonte ancorata, output.
% Tipo: schema concettuale.
%
% \begin{figure}[h]
%     \centering
%     \includegraphics[width=0.95\textwidth]{fig02_grounding_three_ways.pdf}
%     \caption[Tre forme di grounding affiancate]{Tre meccanismi di grounding adottati nella suite: Web grounding (Vera), RAG (CooPolicy) e Text-to-SQL (AllyCare, SysAid).}
%     \label{fig:grounding_three_ways}
% \end{figure}

\subsection{Web grounding}
\label{subsec:bg-web-grounding}
% TODO — Integrazione nativa Google Search in Vertex AI. Quando
% e' appropriato (fatti pubblici, attualita') e quando NO (dati privati).

\subsection{Retrieval-Augmented Generation}
\label{subsec:bg-rag}
% TODO — Pipeline canonica \cite{lewis2020rag}: chunking, embedding,
% vector store, retrieval, generation. Hybrid search dense+sparse,
% re-ranking, query expansion \cite{ma2023query_expansion},
% HyDE \cite{gao2022hyde}. Vertex AI RAG Engine come realizzazione managed.
% Distinzione cruciale (riusata in Cap. 3): RAG corpus per-utente
% vs RAG corpus condiviso.

% [FIGURA 2.2] Pipeline RAG canonica: documento -> chunk -> embedding ->
% indice -> retrieval -> generation con contesto.
%
% \begin{figure}[h]
%     \centering
%     \includegraphics[width=0.9\textwidth]{fig02_rag_pipeline.pdf}
%     \caption[Pipeline RAG canonica]{Pipeline canonica di Retrieval-Augmented Generation: dall'ingestione dei documenti alla generazione vincolata al contesto recuperato.}
%     \label{fig:rag_pipeline}
% \end{figure}

\subsection{Text-to-SQL}
\label{subsec:bg-text2sql}
% TODO — Traduzione NL -> SQL. Challenge: schema awareness, dialetti,
% JOIN multi-tabella, aggregazioni temporali. Benchmark: Spider
% \cite{yu2018spider}, BIRD \cite{li2023bird}. Metriche: execution
% accuracy, exact match. Approcci con LLM: schema injection, few-shot,
% decomposizione \cite{pourreza2023dinsql}.

% =====================================================================
\section{Valutazione automatica di sistemi RAG: il framework RAGAS}
\label{sec:bg-ragas}
% =====================================================================
% TODO — RAGAS \cite{es2023ragas} come insieme di metriche LLM-as-judge:
% faithfulness, answer relevancy, context precision, context recall.
% Limiti del paradigma LLM-as-judge \cite{zheng2023llmjudge}.

% =====================================================================
\section{Agenti LLM e architetture multi-agente}
\label{sec:bg-multiagent}
% =====================================================================
% TODO — Paradigma ReAct \cite{yao2022react} (Thought -> Action -> Observation),
% memoria short/long-term. Survey LLM multi-agente \cite{guo2024multiagent}:
% architetture cooperative, competitive, hierarchical.
% Orion = hierarchical con supervisor agent.
% Pattern AutoGen \cite{wu2023autogen}: GroupChatManager.
% Mixture-of-Agents \cite{wang2024moa} come riferimento terminologicamente
% piu' corretto del classico MoE \cite{jiang2024mixtral} (interno al modello).

% [FIGURA 2.3] Tassonomia architetture multi-agent (cooperative,
% hierarchical, competitive) con esempi di pattern.
%
% \begin{figure}[h]
%     \centering
%     \includegraphics[width=0.85\textwidth]{fig02_multiagent_taxonomy.pdf}
%     \caption[Tassonomia delle architetture multi-agent LLM]{Le tre macro-categorie di architetture multi-agent secondo \cite{guo2024multiagent}, con il posizionamento di Orion nel ramo \emph{hierarchical}.}
%     \label{fig:multiagent_taxonomy}
% \end{figure}

% [FIGURA 2.4] Pattern ReAct: ciclo Thought -> Action -> Observation.
%
% \begin{figure}[h]
%     \centering
%     \includegraphics[width=0.7\textwidth]{fig02_react_loop.pdf}
%     \caption[Ciclo ReAct]{Il ciclo Thought-Action-Observation del pattern ReAct \cite{yao2022react}, alla base degli agenti della suite e del router di Orion.}
%     \label{fig:react_loop}
% \end{figure}

% =====================================================================
\section{LLM routing e ottimizzazione costo--qualita'}
\label{sec:bg-routing}
% =====================================================================
% TODO — Cost-quality trade-off. FrugalGPT \cite{chen2023frugalgpt}:
% LLM cascade, early-exit basato su confidenza. RouteLLM \cite{ong2024routellm}:
% formalizzazione del routing come classificazione strong/weak,
% confronto Matrix Factorization / BERT classifier / Causal LLM /
% SW Ranking. Metodologia di valutazione MMLU/MT-Bench.
%
% Posizionare il router di Orion nel quadro: e' un LLM-based router
% (Gemini Flash) con output strutturato, NON un modello appreso.

% =====================================================================
\section{Stima dell'incertezza e confidence calibration}
\label{sec:bg-confidence}
% =====================================================================
% TODO — Capacita' degli LLM di sapere di non sapere
% \cite{kadavath2022knowsknows}. Self-consistency \cite{wang2022selfconsistency}
% come proxy di confidenza (costoso). Active learning \cite{settles2012active}
% come fondamento del meccanismo di clarification di Orion: il sistema
% interroga l'oracolo (utente) solo quando l'informazione aggiuntiva
% ha massimo valore.

% =====================================================================
\section{Sicurezza, privacy, compliance}
\label{sec:bg-security}
% =====================================================================
% TODO — Model Armor: filtri Responsible AI e difese contro prompt
% injection \cite{perez2022promptinjection}. Cloud DLP: PII con info types
% standard e custom (Codice Fiscale, Partita IVA, IBAN per il contesto
% italiano).

% =====================================================================
\section{Identity federation}
\label{sec:bg-identity}
% =====================================================================
% TODO — OAuth 2.0 \cite{rfc6749oauth2}, OpenID Connect \cite{oidc1_0core},
% JWT \cite{rfc7519jwt}. Microsoft Entra ID e MSAL come implementazione
% adottata in Coop.

% =====================================================================
\section{Infrastruttura cloud}
\label{sec:bg-cloud}
% =====================================================================
% TODO — Cloud Run come piattaforma serverless container.
% Cloud SQL + Private Service Connect. Vertex AI: catalogo modelli +
% RAG Engine + Grounding. Terraform come Infrastructure-as-Code.
% Citare la documentazione ufficiale: \cite{gcp_cloudrun},
% \cite{gcp_vertexai}, \cite{terraform_docs}.
